% ====================================================================
% notation.tex — Ch3(Si·혼합) 기호표 확정안 (W2 초안 — 실측·데이터 강조축)
%   계승 2단 규약(ch3v22_notation 승계): 계승(Ch1·2 정의 그대로 — 재정의 없음) + 신규(본 장 도입).
%   신규 기호는 도입 절을 표에 명시(브리핑 규칙 3). 라벨 네임스페이스 = sec:si-*/eq:si-*/eq:blend-*.
% ====================================================================
\subsection*{기호 --- 계승과 신규(확정안)}
본 장의 기호 규약은 두 단이다. \textbf{계승}(Chapter 1$\cdot$2 정의 그대로 --- \emph{재정의
없음}): 전하 보존 반전~\eqref{eq:sm-mc-balance} 과 그 요동 양성 유일근 논증(\S\ref{sec:sm-mc})$\cdot$평형
진행률 logistic~\eqref{eq:logisticsolve}$\cdot$폭 $w_j=n_jRT/F$$\cdot$전이 용량 $Q_j$$\cdot$방향 부호
$\sigma_d$$\cdot$정규용액 상호작용 $\Omega_j$~\eqref{eq:mu}$\cdot$히스 gap 골격~\eqref{eq:dUhys}$\cdot$관측
평형(개방회로) 전위 $U_\oc(\bar x,T)$~\eqref{eq:sm-mc-balance}. 이들은 흑연$\cdot$LCO 에서 세운 것을
글자까지 그대로 쓰며, 블렌드에서 바뀌는 것은 이 골격에 들어가는 \emph{성분}(host 집합$\cdot$용량 몫)이지
골격 자체가 아니다. 계승 기호의 Si 재해석 한계(격자기체 $\xi$ 의 의미 변경)는 생존 지도
\S\ref{sec:si-survival} 가 노드별로 규정한다.

\textbf{신규}(본 장 도입 --- 표~\ref{tab:si-notation} 가 도입 절의 정의 기준):

\begin{table}[h]
\centering\footnotesize
\renewcommand{\arraystretch}{1.3}
\caption{Chapter 3 신규 기호 --- 도입 절과 정의. 위첨자 host 는 두 활물질 $\{\mathrm{gr},\mathrm{Si}\}$ 를
가르는 표지이고, 나머지는 케이스 실측 파라미터(\S\ref{sec:si-cases})와 응력 결합(\S\ref{sec:si-mech})에서
도입된다. 케이스별 파라미터 셋과 블렌드 합성 기호의 \emph{수치}는 해당 절 표가 정본이다.}
\label{tab:si-notation}
\begin{tabular}{@{}l p{0.50\textwidth} c@{}}
\toprule
기호 & 의미 & 도입 절 \\
\midrule
$f_\mathrm{Si}$        & Si 용량 몫 $Q_\mathrm{Si}/Q$(혼합 음극; 검토 범위 $0$--$0.3$) & \S\ref{sec:si-blend} \\
$Q_\mathrm{gr},Q_\mathrm{Si}$ & host별 총용량, $Q=Q_\mathrm{gr}+Q_\mathrm{Si}$        & \S\ref{sec:si-blend} \\
$Q_j^{\mathrm{host}}$  & host 내 전이 $j$ 의 용량(host $\in\{\mathrm{gr},\mathrm{Si}\}$) & \S\ref{sec:si-blend} \\
$\xi_{\eq,j}^{\mathrm{host}}$ & host 내 전이 $j$ 의 평형 진행률(공유 $U_\oc$ 에서 평가)   & \S\ref{sec:si-blend} \\
케이스 표지            & 원소 Si $/$ SiO$_x$ $/$ Si--C 복합(활물질 계열)          & \S\ref{sec:si-cases} \\
$\eta_\mathrm{ICE}$    & 1차 쿨롱효율(1차 비가역 지표, 실측 \%)                    & \S\ref{sec:si-cases} \\
$\partial U/\partial T$& 엔트로피 계수(측정 열역학 서명, $\mu$V/K)                & \S\ref{sec:si-cases} \\
$V_\mathrm{Li}$        & host 내 Li 부분몰 부피(응력--조성 결합)                   & \S\ref{sec:si-mech} \\
$\sigma_h$             & 정수압(평균) 응력 $\tfrac13\operatorname{tr}\sigma$      & \S\ref{sec:si-mech} \\
$\Delta U_\sigma$      & 응력 유발 평형 중심 이동 $V_\mathrm{Li}\sigma_h/(sF)$     & \S\ref{sec:si-mech} \\
\bottomrule
\end{tabular}
\end{table}

두 표기 관행을 덧붙인다. (i) host 위첨자는 \emph{자리 클래스의 소속}만 가르는 표지이고 계승 아래첨자
$j$(전이)$\cdot w_j\cdot Q_j$ 의 의미는 불변이다 --- 곧 $Q_j^{\mathrm{host}}$ 는 ``host 의 전이 $j$'' 이지 새
물리량이 아니다. (ii) 관측 평형 전위는 Ch1 의 $U_\oc$ 를 그대로 쓰되, 블렌드에서는 두 host 가
\emph{같은} $U_\oc$ 를 공유한다는 점(\S\ref{sec:si-blend} 의 공유 저장조)만 새로우며 기호는 재정의하지
않는다. Si 계열에서 아래첨자 $j$ 가 ``진짜 상전이'' 가 아니라 유효 곡선 기술 성분일 수 있다는 재해석
(\S\ref{sec:si-survival} N5)은 의미론적 각주이지 기호 변경이 아니다.

% 자산(v1.0.22 R5 신규): [V22-R5-01 Ch3 기호표 확정안(계승 2단 + 신규 10 기호 도입절 표 tab:si-notation)]
